% Source: https://artofproblemsolving.com/community/c3h1725653_2018_mathcounts_national_target_2
% Stable identity: aops-topic-1725653-q01
Let $T$ be the set of $n$ for which the $n$th number is $2n$; for the others it is $\dfrac {n}{2}$. The sum must be $20\cdot12=240$:
\[
\sum_{n\in T}2n+\sum_{n\notin T}\dfrac {n}{2}=\dfrac{1}{2}\cdot210+\dfrac{3}{2}\sum_{n\in T}n=240\quad\Longrightarrow\quad\sum_{n\in T}n=90 .
\]
The product is $P=20!\cdot2^{|T|}\cdot\left(\dfrac{1}{2}\right)^{20-|T|}$, so $\dfrac{P}{20!}=2^{2|T|-20}$, which is largest when $T$ has as many elements as possible.

Since $1+2+\cdots+13=91>90$, any $13$ distinct numbers from $1$ to $20$ add up to more than $90$, so $|T|\le12$. Twelve is possible: $T=\{1,2,\dots,10,15,20\}$ has sum $55+35=90$. So the greatest value is $2^{24-20}=\boxed{16}$.
